\documentclass[11pt]{article}
\usepackage{amsmath,amsfonts}
\usepackage{graphicx}
%\usepackage{xcolor}
%\definecolor{gray}{rgb}{.75,.75,.75}
%\usepackage[landscape]{geometry}
%\usepackage{hyperref}
%\usepackage[bookmarks=true, pdfborder={0 0 .5}, urlbordercolor=gray]{hyperref}

\renewcommand{\thefootnote}{\fnsymbol{footnote}}


\addtolength{\topmargin}{-.6in}
\addtolength{\textheight}{1.2in}
\addtolength{\oddsidemargin}{-.2in}
\addtolength{\textwidth}{.4in}
\pagestyle{empty}

\newcommand{\ds}{\displaystyle}
\newcommand{\ts}{\textstyle}

\newcommand{\Z}{\mathbb{Z}}
\newcommand{\R}{\mathbb{R}}
\newcommand{\C}{\mathbb{C}}
\newcommand{\EE}{\mathcal{E}}
\newcommand{\tZ}{\tilde{\Z}}

\newcommand{\Sym}{\mathrm{Sym}}

\renewcommand{\circle}[1]{{\bigcirc\hspace{-.75em}#1}}

\begin{document}

\footnote[0]{Notes by Peter Magyar \ {\tt magyar@math.msu.edu}}
\vspace{-1em}

\centerline{\bf Math 133\hspace{1.5em} \hfill{\Large\bf
Integral Test}\hfill Stewart \S11.3}
\vspace{1em}

\noindent 
{\bf Series and integrals.} 
Our goal for infinite series is to express a complicated quantity
as an infinite sum of simple terms, so that finite partial sums 
approximate the quantity as accurately as we like.
We will not have significant tools to achieve this until \S11.9.

For now (\S11.3-11.7), we concentrate on a less specific question: 
when does a given series converge to {\it some} finite value, 
rather than diverge to infinity or an oscillating value?
For example, we have seen the $n$th Term Vanishing Test: 
the series must diverge if the the terms do not approach zero, $\lim_{n\to\infty}a_n\neq 0$.
A more subtle and powerful convergence test comes
from comparing the sum of a series to the area under a curve
$y=f(x)$ passing through each point $(n,a_n)$.

\begin{quote}
{\it Integral Test:} Suppose the function $f(x)$ is continuous,
positive, and decreasing on the interval $x\in[1,\infty)$,
and that $a_n=f(n)$. We compare the improper integral  $\int_1^\infty f(x)\,dx$
with the infinite series $\sum_{n=1}^\infty a_n$.

\begin{itemize}
\item If $\int_1^\infty f(x)\,dx$ diverges, 
then $\sum_{n=1}^\infty a_n$ also diverges.
\item If $\int_1^\infty f(x)\,dx$ converges, 
then $\sum_{n=1}^\infty a_n$ also converges.
\end{itemize}
\end{quote}
\medskip

\noindent
{\bf Divergent case.} 
Consider $\sum_{n=1}^\infty a_n=\sum_{n=1}^\infty \frac1n$.
Then the function $f(x)=\frac1x$ has $a_n=f(n)$,
and for $x\in[1,\infty)$, this function is: 
\begin{itemize}
\item continuous, since its only vertical asymptote is $x=0$, outside $x\in[1,\infty)$\,;
\item positive, since $x\geq 1$ implies $\frac1x>0$\,;
\item decreasing, since its derivative is negative,\footnote{
See \S3.3 Derivatives and Graphs, in the Math 132 Lecture Notes} $f'(x)=-\frac1{x^2}<0$.
\end{itemize}
Thus, we can apply the Integral Test to compare the infinite series with the improper integral  $\int_1^x f(x)\,dx$. We compute:
$$\textstyle
\int_1^\infty \frac1x \,dx 
\ =\ \lim\limits_{n\to\infty} \int_1^n \frac1x \,dx 
\ =\ 
 \left.\lim\limits_{n\to\infty}\ln(x)\right|_{x=1}^{x=n}
\ =\ 
\lim\limits_{n\to\infty} \ln(n)-\ln(1) \ =\ \infty.
$$
Since the integral diverges, the series $\sum_{n=1}^\infty a_n$ must also diverge.\footnote{
Elementary proof of divergence: $\textstyle \sum_{n=1}^\infty \frac1n = 1 +\underbrace{ \tfrac12}_{\geq1/2}
+\underbrace{\tfrac13+\tfrac14}_{\geq1/2}
+\underbrace{\tfrac15+\tfrac16+\tfrac17+\tfrac18}_{\geq 1/2}+\cdots = \infty$.
}

The Integral Test is best understood geometrically.
The value of the series is the total area under the bar-graph 
of $\{a_n\}$, where we draw the bar at height $a_n$ 
above the interval $x\in[n,n{+}1]$.
\\[1em]
\centerline{
\includegraphics[width=3.5in]
{11-3-Integral-Divergence.png} 
}
\vspace{-1em}

\noindent
The integral is the area of the region under
$y=f(x)$ and above $x\in[1,\infty)$.
But the bar graph completely contains the integral region, so:
$$
\sum_{n=1}^\infty a_n \ \geq \ \int_1^\infty f(x)\,dx\,.
$$ 
Since the integral diverges to infinity,
getting larger and larger with no bound as we add area on the right,
the series must also diverge as we add more terms.
\\[1em]
\noindent
{\bf Convergent case.} Now consider 
$\sum_{n=1}^\infty a_n=\sum_{n=1}^\infty \frac1{n^2}$
which has  $a_n=f(n)$ for $f(x)=\frac1{x^2}$, which is again
a continuous, positive, decreasing function for $x\in[1,\infty)$.
Thus we can apply the Integral Test to compare the infinte series 
with the improper integral $\int_1^x f(x)\,dx$:
$$\textstyle
\int_1^\infty \frac1{x^2} \,dx 
\ =\ \lim\limits_{n\to\infty} \int_1^n \frac1{x^2} \,dx 
\ =\ 
 \left.\lim\limits_{n\to\infty}-\frac1x\right|_{x=1}^{x=n}
\ =\ 
\lim\limits_{n\to\infty} -\frac1n-(-\frac11) \ =\ 1.
$$
Since the integral converges, the series $\sum_{n=1}^\infty a_n$ must also converge.

Again, we can understand this geometrically. 
The value of the series is the total area under the bar-graph 
of $\{a_n\}$, but this time we draw the bar at height $a_n$ 
above the interval $x\in[n{-}1,n]$, shifted left from our previous picture, 
so the bar graph lies under $y=f(x)$.
For this series, the heights are $a_n=\frac1{n^2}$, much lower than
the previous $a_n=\frac1n$, so it has a better
chance of converging to a finite area.
\\[1em]
\centerline{
\includegraphics[width=3.5in]
{11-3-Integral-Convergence.png} 
}
\vspace{-1em}

\noindent
Clearly, the part of the bar-graph after $a_2$ is contained in the integral region
under $y=f(x)$ and above $x\in[1,\infty)$. Thus:
$$
\sum_{n=2}^\infty a_n \ \leq \ \int_1^\infty\!\! f(x)\,dx\,, 
$$ 
$$
\sum_{n=1}^\infty a_n \ =\ a_1+\sum_{n=2}^\infty a_n  \ \leq \ a_1 + \int_1^\infty\!\! f(x)\,dx \ =\ 2\,.
$$ 
So $\sum_{n=1}^\infty \frac1{n^2}$ is not only finite, it is less than 2.\footnote{
In fact, Euler computed this limit as $\frac{\ \pi^2}{6}\approx 1.64$.
Look up the Basel Problem.}
\\[1em]
{\bf Examples.} \begin{itemize}

\item {\it Standard p-series.} 
The above reasoning is easily generalized to show:

$$\boxed{
\ \sum_{n=1}^\infty \frac1{n^p}
\ \ \text{converges if $p>1$, diverges if $p\leq1$.}\ 
}$$
\\[-1.5em]

\item  Determine the convergence of:
$$
\sum\limits_{n=1}^\infty a_n \ =\ \sum\limits_{n=1}^\infty \frac1{(2n-5)(2n-7)}\,. 
$$
In this case, the function $f(x)=\frac1{(2x-5)(2x-7)}$
has vertical asymptotes at $x=\frac52,\frac72$:
it is not continuous, or positive, or decreasing on $x\in[1,\infty)$,
so the Integral Test does not immediately apply.

However, to the right of the asympotes, for $x\in[4,\infty)$,
the function is continuous; it is positive, 
since $2x-5>0$ and $2x-7>0$ for $x\geq 4$;
and it is decreasing, since the derivative
$f'(x) =  -\frac{8 (x-3)}{(5-2 x)^2 (7-2 x)^2} < 0$
for $x\geq 4$. Further, we have:
$$\begin{array}{rcl}
\int_4^\infty \frac1{(2x-5)(2x-7)}\, dx
&=&
\int_4^\infty\!\! -\frac{\frac12}{2 x-5} +  \frac{\frac12}{2 x-7}\  dx
\\[1em]
&=& 
\lim_{n\to\infty}
\left.
\tfrac14 \ln\!\left(\frac{2x{-}7}{2x{-}5}\right)
\right|_{x=4}^n
\ =\
0 - \ln(\frac13) \ < \ \infty.
\end{array}$$
since $\lim_{n\to\infty} \ln\!\left(\frac{2n{-}7}{2n{-}5}\right)
=\ln(1)=0$.

Slightly generalizing the reasoning of the convergent case of the Integral Test above,
we find that $\sum\limits_{n=5}^\infty a_n \leq \int_4^\infty f(x)\,dx$.
Thus, we have:
$$\begin{array}{rcl}
\sum\limits_{n=1}^\infty a_n
&=&
a_1+a_2+a_3+a_4+ \sum\limits_{n=5}^\infty a_n
\\[1.5em]
&<& a_1+a_2+a_3+a_4+ \int_4^\infty f(x)\,dx \ <\ \infty.
\end{array}$$

\end{itemize}

\end{document}

%%%%%%%%%%%  Picture  %%%%%%%%%%
\\[1em]
\centerline{
%\includegraphics[width=2in]
{1-8-Discont-iv.png}
\qquad \qquad 
%\includegraphics[width=2in]
{1-8-Discont-v.png} 
}
\vspace{0em}

\noindent
%%%%%%%%%%%%%%%%%%%%%%%%%

%%%%%%%%%%%  Table  %%%%%%%%%%
In fact, we get the following derivatives:
$$\begin{array}{|c||c|c|c|c|c|c|}
\hline &&&&&& \\[-.9em] 
f(x) & \sin(x) & \cos(x) & \tan(x)  & \sec(x) & \csc(x) & \cot(x)
\\[.2em] \hline &&&&&& \\[-.9em]
f'(x) &\cos(x) & -\sin(x) & \sec^2(x) & \tan(x)\sec(x) & -\cot(x)\csc(x) & -\csc^2(x)
\\[.1em] \hline
\end{array}$$
\\[1em]
%%%%%%%%%%%%%%%%%%%%%%%%%%



















